% report.tex -- CS601T Deep Learning, Programming Assignment 5 (Group 11): CNNs.
% Overleaf: upload this folder (report.tex + figures/), compile with pdfLaTeX.

\documentclass[11pt,a4paper]{article}
\usepackage[utf8]{inputenc}
\usepackage[T1]{fontenc}
\usepackage{lmodern}
\usepackage[margin=1in]{geometry}
\usepackage{amsmath,amssymb}
\usepackage{booktabs}
\usepackage{graphicx}
\usepackage[font=small,labelfont=bf,skip=4pt]{caption}
\usepackage{subcaption}
\usepackage{float}
\usepackage[table]{xcolor}
\usepackage{enumitem}
\usepackage{microtype}
\usepackage{array}
\usepackage{fancyhdr}
\usepackage[colorlinks=true,linkcolor=blue!55!black,citecolor=blue!55!black,urlcolor=blue!55!black]{hyperref}

\graphicspath{{figures/}}

\definecolor{selrow}{HTML}{E3ECF7}
\definecolor{inferencebg}{HTML}{F4F6F8}
\definecolor{inferencebar}{HTML}{4A6FA5}
\newcommand{\infer}[1]{%
  \par\vspace{2pt}\noindent
  \colorbox{inferencebg}{\begin{minipage}{\dimexpr\linewidth-2\fboxsep\relax}
    \vspace{2pt}\small\textbf{\textcolor{inferencebar}{What this shows}} --- #1\vspace{2pt}
  \end{minipage}}\par\vspace{4pt}}
\newcommand{\code}[1]{\texttt{#1}}

\setlength{\parskip}{3pt}\setlength{\parindent}{0pt}\renewcommand{\arraystretch}{1.12}
\pagestyle{fancy}\fancyhf{}
\fancyhead[L]{\small\itshape CS601T: Deep Learning --- Assignment 5}
\fancyhead[R]{\small\itshape Group 11}
\fancyfoot[C]{\small\itshape\thepage}
\renewcommand{\headrulewidth}{0.4pt}

\begin{document}

\begin{titlepage}
\centering
\vspace*{1em}
{\large\scshape CS601T: Deep Learning\par}\vspace{0.2em}{\normalsize Programming Assignment 5\par}
\vspace{1.2em}\rule{\linewidth}{1.0pt}\\[0.6em]
{\LARGE\bfseries Convolutional Neural Networks:\\[0.2em] Building Blocks, Architectures and Visualisation\par}
\vspace{0.6em}\rule{\linewidth}{1.0pt}\vspace{2.4em}
{\large\bfseries Group 11\par}\vspace{0.8em}
\begin{tabular}{@{}ll@{}}
\toprule\textbf{Roll Number} & \textbf{Name} \\ \midrule
CS24BT055 & Himanshu Nagwanshi \\
CS24BT060 & Anshuman Rahul \\
CS24BT061 & Sidharth B \\
\bottomrule
\end{tabular}
\vspace{2.4em}
\begin{tabular}{@{}ll@{}}
\toprule\textbf{Course} & CS601T: Deep Learning \\ \midrule
\textbf{Semester} & August--November 2026 \\
\textbf{Institute} & Indian Institute of Technology Dharwad \\
\bottomrule
\end{tabular}
\vfill
{\small Code: \code{Group11\_Assignment5\_code}\par}
\vspace{1em}
\end{titlepage}

\begin{abstract}
\noindent
Part~A builds a convolutional network for a five-class MNIST subset (digits 0, 4, 5, 6, 7) one
block at a time: a single Kaiming-initialised $5\times5$ filter applied by explicit traversal, a
convolution--ReLU--max-pool layer class stacked twice, and the full network trained with SGD.  The
trained network reaches 99.42\,\% test accuracy (22 errors in
3795).  Part~B trains three CNNs from scratch on a five-class Caltech-101 subset
($224\times224$ colour images, 250 training images).  All three fit their training set almost
perfectly (97.2--99.6\,\%) while reaching 74.0--78.0\,\% validation accuracy; the
5-layer Arch-2 is selected on validation and scores 68.0\,\% on the test set; rerunning the
whole procedure with five seeds finds no evidence that one architecture is better (70.0--71.6\,\% mean
test accuracy).  An
ImageNet-pretrained VGG19 with only its classification layer retrained scores
100.0\,\% on the same test set.  Every network is also examined from the inside:
filters, feature maps, the input patches that maximally activate a neuron, guided
backpropagation and Grad-CAM.
\end{abstract}

\tableofcontents
\newpage

% =====================================================================================
\section{Data and Common Method}\label{sec:method}

\subsection{Data}
\begin{table}[H]\centering\small
\begin{tabular}{@{}l>{\raggedright\arraybackslash}p{3.6cm}lrrr@{}}\toprule
dataset & classes & image & train & val & test \\ \midrule
MNIST subset & 0, 4, 5, 6, 7 & $28\times28$ grey & 11385 & 3795 & 3795 \\
Caltech-101 subset & bonsai, buddha, car\_side, scorpion, starfish & $224\times224\times3$ (resized)
  & 250 & 50 & 100 \\ \bottomrule
\end{tabular}
\caption{Group 11 datasets.  Both are balanced across classes (MNIST 2277/759/759 per class,
Caltech 50/10/20), so chance accuracy is 20\,\%.  No file appears in more than one split (checked
by content hash).}\label{tab:data}
\end{table}
MNIST pixels are scaled to $[0,1]$.  Caltech images vary in size (124 to 300 pixels per side) and
are resized to $224\times224$ without preserving aspect ratio, as the sheet asks.  Of the 400
Caltech images, all 80 car\_side images are greyscale, as are 4 buddha and 4 starfish images (one
starfish file is stored as RGB but has three identical channels); every greyscale image is
replicated to three channels.  Because car\_side is the only class that is
entirely greyscale, a network can partly recognise it from the absence of colour; we return to this
in Section~\ref{sec:limits}.  For the networks trained from scratch the Caltech pixels are
normalised per channel with the mean and standard deviation of the training split only.  VGG19
receives the same resized pixels through its own \code{vgg19.preprocess\_input}.

The images used for every visualisation are fixed by rule before training: the first file of each
class in the training split, and for the single-image questions (B1c, B2c) the first bonsai image,
so that the colour feature maps have colour to show.

\subsection{Layers, sizes and initialisation}
Every convolution is followed by ReLU (as the lecture notes prescribe for convolution outputs), except
A1's single filter, whose raw response the sheet asks for.  The
output side of a convolution or pooling layer with input width $W_1$, kernel $F$, padding $P$ and
stride $S$ is
\begin{equation}\label{eq:size}
W_2 = \left\lfloor \frac{W_1 - F + 2P}{S} \right\rfloor + 1 ,
\end{equation}
with the depth equal to the number of filters.  The floor matters in Part~B: $(224-11)/4 = 53.25$,
so the first Caltech layer is $54\times54$ and its last input column is never covered by a filter.
Every layer size reported below was predicted with Eq.~\eqref{eq:size} and asserted against the
tensor the code actually produced.  Convolution kernels are Kaiming-normal initialised~\cite{he}
($\mathcal N(0, 2/\text{fan-in})$), fully connected layers Xavier-normal~\cite{glorot}, biases zero.  The output
layer is a 5-way softmax trained with cross-entropy.

\subsection{Training protocol}\label{sec:protocol}
The sheet fixes the architectures and the initialisation and asks for stochastic gradient descent,
but does not fix its settings, the stopping rule or the learning rate.  One protocol is used for
every trained model, so that all comparisons share it:
\begin{itemize}[leftmargin=*,itemsep=1pt]
\item \textbf{Optimiser.} Mini-batch stochastic gradient descent with momentum $0.9$, no weight
  decay, no learning-rate schedule.  Momentum keeps the method SGD (gradients from
  backpropagation on random mini-batches); it only smooths the update direction.  Batch size 32 for MNIST and 16 for every Caltech model
  (32 would give only 8 updates per epoch on 250 images).
\item \textbf{Stopping.} Early stopping on the validation loss: training stops once the validation
  loss has gone 15 epochs without a new minimum, and the weights of the best epoch are restored.
  A run whose loss becomes infinite or NaN is stopped at once, and runs are capped at 200 epochs
  (MNIST) and 500 epochs (Caltech, VGG19).
\item \textbf{Learning rate.} Every model runs the same grid $\{0.001, 0.003, 0.01, 0.03, 0.1\}$ from
  identical initial weights and keeps the rate with the lowest early-stopped validation loss.  If
  that rate is at an end of the grid, one further step ($0.0003$ or $0.3$) is tried once.  Each
  model therefore gets its own rate, chosen by an identical procedure with an identical budget.  A
  single shared rate would instead favour whichever architecture happens to suit it.
\item \textbf{Fairness of the Part~B comparison.} The three Caltech architectures share the data,
  preprocessing, batch size, optimiser, stopping rule, learning-rate grid, the initial weights of
  their two common convolution layers, and the order in which training images are presented.  The
  VGG19 head uses the same batch size, optimiser, grid and stopping rule.
\item \textbf{Use of the splits.} The validation set chooses the stopping epoch, the learning rate
  and (B1b) the architecture.  The test set is evaluated only after all of these choices, and only
  for the models the sheet asks about.  The one addition is the seed study of
  Section~\ref{sec:seeds}, which records the test accuracy of all three Caltech architectures after
  each seed's choices are made, for analysis only.
\end{itemize}
Every run of every grid is reported (Tables~\ref{tab:a3grid}, \ref{tab:b1grid} and \ref{tab:b2grid}).

\subsection{Tracing a neuron back to its input patch}\label{sec:rf}
The patch that ``causes a neuron to fire''~\cite{zeiler} is found as in the lecture notes: starting from the
neuron, step back one layer at a time, widening the window by that layer's kernel and stride, until
the input is reached.  For a layer with kernel $k$, stride $s$ and padding $p$, an output window of
rows $[r_0, r_1]$ comes from input rows
\begin{equation}\label{eq:trace}
[\,r_0 s - p,\;\; r_1 s - p + k - 1\,],
\end{equation}
applied to every convolution and pooling layer from the neuron down to the image (columns
likewise).  Whole pooling windows are kept, not only the position that won the maximum.  The
resulting box is the neuron's receptive field; its side grows as $\text{rf} \leftarrow
\text{rf} + (k-1)\,j$ with $j$ the product of the strides below.  The code checks this box against
the exact set of input pixels with a non-zero gradient for randomly chosen neurons of every
network, and the two agree.  For this check ReLU is replaced by the identity and max pooling by
average pooling, so that every pixel the neuron can depend on receives gradient; with them, the
gradient reaches only a subset of the box.  ``Last convolution layer'' always means its output after ReLU and
before pooling.  Neuron positions are given 1-based as (row, column) in feature map $m$.

% =====================================================================================
\section{Part A: CNN Building Blocks on MNIST}\label{sec:a}

\subsection{A1 -- One filter, traversed by hand}
One $5\times5$ filter is drawn from $\mathcal N(0, 2/25)$ (Kaiming, fan-in 25) and slid over each
of the five images by explicit loops, stride 1, padding 0.  The traversal computes
$s_{ij} = \sum_{a,b} I_{i+a,\,j+b}\,K_{ab}$, the notes' image formula, which is what PyTorch and
Keras compute as well; the code checks it against \code{conv2d} to $10^{-9}$.
Eq.~\eqref{eq:size} gives $(28-5+0)/1+1 = 24$, and all five obtained maps are $24\times24$.
No activation is applied, so the maps are signed.

\begin{table}[H]\centering\small\begin{tabular}{@{}rrrrr@{}}\toprule
0.010 & 0.385 & 0.346 & -0.144 & -0.084 \\
-0.149 & 0.161 & -0.016 & 0.211 & -0.523 \\
0.443 & -0.027 & 0.192 & -0.039 & -0.107 \\
0.131 & 0.233 & -0.057 & -0.043 & 0.194 \\
-0.246 & -0.428 & 0.112 & -0.190 & -0.543 \\
\bottomrule\end{tabular}
\caption{A1 filter values (also shown, rounded, in Figure~\ref{fig:a1}).}\label{tab:a1filter}
\end{table}

\begin{figure}[H]\centering
\begin{minipage}[c]{0.36\linewidth}\centering\includegraphics[width=\linewidth]{A/A1_filter.png}\end{minipage}\hfill
\begin{minipage}[c]{0.40\linewidth}\centering\includegraphics[width=\linewidth]{A/A1_feature_maps.png}\end{minipage}
\infer{A random filter is already a feature detector: it responds with opposite signs on the two
sides of a stroke, so each map traces the digit as a pair of bright (yellow) and dark (purple)
bands, with near-zero response (teal) on the empty background.  Nothing here is learned: the structure comes from the images,
and the filter only decides which local contrast pattern each output value measures.}
\caption{A1.  Left: the filter.  Right: the five images and their $24\times24$ feature maps
(one shared viridis scale, yellow most positive, purple most negative).  As in the notes, filters
are drawn in grey scale (white most positive, black most negative, mid-grey zero) and feature maps
in viridis.}\label{fig:a1}
\end{figure}

\subsection{A2 -- A convolution--ReLU--pool class, stacked twice}
\code{ConvReLUPool} (in \code{layers.py}) holds its own Kaiming-initialised weights and computes
convolution, ReLU and $2\times2$ max pooling with stride 2.  The convolution itself uses PyTorch's
\code{conv2d}, the operation A1 verifies against the explicit loop.  Two instances are stacked: 6 filters
of $5\times5\times1$, then 16 filters of $5\times5\times6$.  The untrained layers are applied to the
A1 images; Table~\ref{tab:a2dims} compares every stage with Eq.~\eqref{eq:size}.  The filters
shown are numbers 1, 3, 5 of layer~1 and 1, 6, 11 of layer~2, the same positions A3b shows
after training.  The program also prints the layer outputs; the full arrays are in the results
JSON.

\begin{table}[H]\centering\small\begin{tabular}{@{}lll@{}}\toprule stage & formula & obtained \\ \midrule
layer1 conv & $6\times24\times24$ & $6\times24\times24$ \\
layer1 pool & $6\times12\times12$ & $6\times12\times12$ \\
layer2 conv & $16\times8\times8$ & $16\times8\times8$ \\
layer2 pool & $16\times4\times4$ & $16\times4\times4$ \\
\bottomrule\end{tabular}
\caption{A2 output sizes (channels $\times$ height $\times$ width), formula vs obtained.}\label{tab:a2dims}
\end{table}

\begin{figure}[H]\centering
\includegraphics[width=0.62\linewidth]{A/A2_filters_l1.png}\\[2pt]
\includegraphics[width=0.85\linewidth]{A/A2_filters_l2.png}
\caption{A2 untrained filters.  Top: three layer-1 filters.  Bottom: three layer-2 filters; each
has one $5\times5$ slice per input channel, drawn as a row of six, as the notes draw VGG's
$3\times3\times3$ filters.}\label{fig:a2filters}
\end{figure}

\begin{figure}[H]\centering
\includegraphics[width=0.74\linewidth]{A/A2_maps_l1.png}\\[4pt]
\includegraphics[width=0.74\linewidth]{A/A2_maps_l2.png}
\infer{After ReLU only one sign of each filter's response survives, so each layer-1 map keeps one
side of the strokes; pooling halves the resolution but keeps the shape legible at $12\times12$.  At
layer~2 the maps no longer look like the digit: each $8\times8$ value is computed from a
$14\times14$ input patch and mixes all six layer-1 maps with random weights.  Untrained, the
responses are spread out: f1 and f6 are non-zero almost everywhere (94\,\% and 92\,\% of values),
while f11 is almost always cut off by ReLU (8\,\%).  After training
(Figure~\ref{fig:a3bmaps}), the same filters fire only on specific parts of the digit, and the
share of non-zero layer-2 values falls from 53\,\% to 19\,\%.}
\caption{A2 feature maps of the chosen filters, after conv+ReLU and after pooling, for the five
images.  The colour scale is shared down each column, so equal colour means equal activation for a
filter across images.}\label{fig:a2maps}
\end{figure}

\subsection{A3 -- The full CNN}
The network, a LeNet-style design~\cite{lecun}, is input $28\times28$ $\rightarrow$ 6@$24\times24$ $\rightarrow$ pool 6@$12\times12$
$\rightarrow$ 16@$8\times8$ $\rightarrow$ pool 16@$4\times4$ $\rightarrow$ flatten 256
$\rightarrow$ 64 (ReLU) $\rightarrow$ 5 (softmax), with 19,345 parameters.  Its two
convolution layers start from exactly the A2 weights, so A3b compares the same filters before and
after training.  Table~\ref{tab:a3grid} gives the learning-rate grid.

\begin{table}[H]\centering\small\begin{tabular}{@{}lrrrrr@{}}\toprule
lr & best ep. & stop ep. & best val loss & val acc (\%) & train acc (\%) \\ \midrule
0.001 & 40 & 55 & 0.0292 & 98.95 & 99.9 \\
0.003 & 20 & 35 & 0.0281 & 99.03 & 100.0 \\
0.01 & 11 & 26 & 0.0214 & 99.21 & 100.0 \\
\rowcolor{selrow}0.03 & 6 & 21 & 0.0190 & 99.34 & 99.9 \\
0.1 & 1 & 16 & 0.6278 & 85.67 & 86.0 \\
\bottomrule\end{tabular}
\caption{A3 learning-rate grid (Section~\ref{sec:protocol}).  The shaded row is selected.  Accuracies
are at the restored best epoch.}\label{tab:a3grid}
\end{table}

\paragraph{A3a -- Accuracy.}  The selected network (learning rate 0.03, best epoch
6) reaches 99.95\,\% on the training set, 99.34\,\% on
validation and 99.42\,\% on test.  The largest cell of the test confusion matrix is a
digit~6 read as 5: 8 of the
22 test errors.  On the same data and splits, the best fully connected network of Assignment~3
($784\rightarrow128\rightarrow64\rightarrow32\rightarrow5$, 110,981 parameters) reached
98.76\,\% test accuracy (47 errors): the CNN makes fewer than half the errors with about a sixth of
the parameters, the gain from small filters shared across all image positions.

\begin{figure}[H]\centering
\begin{subfigure}{0.32\linewidth}\includegraphics[width=\linewidth]{A/A3a_cm_train.png}\caption{train}\end{subfigure}\hfill
\begin{subfigure}{0.32\linewidth}\includegraphics[width=\linewidth]{A/A3a_cm_val.png}\caption{validation}\end{subfigure}\hfill
\begin{subfigure}{0.32\linewidth}\includegraphics[width=\linewidth]{A/A3a_cm_test.png}\caption{test}\end{subfigure}
\caption{A3a confusion matrices (counts and row percentages).}\label{fig:a3cm}
\end{figure}

\begin{figure}[H]\centering
\includegraphics[width=0.85\linewidth]{A/A3a_curves.png}
\infer{The validation loss reaches its minimum at epoch 6 and then drifts up while the
training loss keeps falling by orders of magnitude; validation accuracy is flat from epoch 2.  The
extra epochs memorise the training set without improving held-out accuracy, which is why the best
epoch, not the last, is restored.}
\caption{A3 training curves of the selected run.  Training values are running averages over each
epoch; validation values are measured after it.}\label{fig:a3curves}
\end{figure}

\paragraph{A3b -- Filters and feature maps after training.}  The sheet asks for 10 filters of the
first convolution layer, which has 6, so all 6 are shown, together with 10 of the 16 layer-2
filters (including the three A2 positions).

\begin{figure}[H]\centering
\includegraphics[width=0.95\linewidth]{A/A3b_filters_l1.png}
\infer{Layer-1 filters can be read directly: filter 1 (a bright middle column next to dark ones)
and filter 5 (a bright right column) respond to vertical and steep diagonal strokes, and filter 3
(a bright top row) to the top edge of a horizontal stroke.}
\caption{A3b: all 6 trained layer-1 filters.}\label{fig:a3bfilters}
\end{figure}

\begin{figure}[H]\centering
\includegraphics[width=0.62\linewidth]{A/A3b_before_after_l1.png}
\infer{Training refines the random filters rather than replacing them: each keeps much of its
initial sign pattern while its contrast grows (filter 3, for example, develops a strong positive
top row).  The relative change $\lVert W_{\text{after}}-W_{\text{before}}\rVert /
\lVert W_{\text{before}}\rVert$ ranges from 0.16 to 1.00 across the six
layer-1 filters and from 0.40 to 1.57 across the ten layer-2 filters shown.
Filter 6, which moved least (relative change 0.16), also gives the weakest maps in
Figure~\ref{fig:a3bmaps}.}
\caption{A3b layer-1 filters at the A2 positions before (A2) and after training.  The layer-2
version is in the code's output (\code{A3b\_before\_after\_l2.png}).}\label{fig:a3bba}
\end{figure}

\begin{figure}[H]\centering
\includegraphics[width=0.66\linewidth]{A/A3b_filters_l2.png}
\infer{Layer-2 filters cannot be read as images: each $5\times5$ slice weights a layer-1 feature
map, not pixels, so their patterns have no visual meaning.  As the lecture notes say, filter images
are interpretable only for the first layer; what layer~2 has learned is shown by its feature maps
(Figure~\ref{fig:a3bmaps}).}
\caption{A3b: 10 of the 16 trained layer-2 filters, one $5\times5$ slice per input channel.  The
three A2 positions (1, 6, 11) come first.}\label{fig:a3bfilters2}
\end{figure}

\begin{figure}[H]\centering
\includegraphics[width=0.85\linewidth]{A/A3b_maps_l1.png}\\[4pt]
\includegraphics[width=0.95\linewidth]{A/A3b_maps_l2.png}
\infer{The trained layer-1 maps are more selective than the untrained ones in
Figure~\ref{fig:a2maps}: individual filters now pick out stroke directions (one fires on the
right-leaning diagonal of the 7 and the 4, another on horizontal bars), and one (filter 2) covers
the whole stroke.  Layer-2 maps are sparse and differ between digits, which is what the classifier
on top of them needs.}
\caption{A3b trained feature maps (conv+ReLU) of the five images: layer 1 (top, all 6 filters) and
layer 2 (bottom, 10 filters).}\label{fig:a3bmaps}
\end{figure}

\paragraph{A3c -- Patches that maximally activate the last convolution layer.}  The last
convolution layer is conv2's $16\times8\times8$ output.  For each image the five neurons chosen are
the five feature maps with the highest peak activation, each at its peak position (taking the five
largest values outright tends to return neighbouring positions of one map, i.e.\ five nearly
identical boxes).  The neurons are chosen per image, so each patch shows what fired most on that
image, not what one fixed neuron prefers; some maps still recur, map 14 among the five for every
digit, and so is map 6 (Table~\ref{tab:a3c}).
Tracing back with Eq.~\eqref{eq:trace}, the neuron at 0-based position $(i,j)$
sees input rows and columns $2i$ to $2i+13$: a $14\times14$ patch.

\begin{table}[H]\centering\small\begin{tabular}{@{}lccccc@{}}\toprule image & 1st & 2nd & 3rd & 4th & 5th \\ \midrule
digit 0 & (6,8,3) 5.9 & (3,5,1) 5.1 & (14,3,5) 4.3 & (5,6,3) 4.0 & (1,1,6) 4.0 \\
digit 4 & (6,5,1) 4.5 & (14,7,1) 4.3 & (16,2,8) 3.5 & (15,6,2) 3.0 & (8,7,7) 2.8 \\
digit 5 & (14,4,7) 4.6 & (8,3,7) 4.5 & (6,2,6) 3.9 & (10,5,7) 3.1 & (16,8,6) 2.9 \\
digit 6 & (6,7,2) 4.1 & (14,8,5) 3.5 & (8,8,4) 3.1 & (13,6,1) 3.0 & (12,2,3) 2.9 \\
digit 7 & (16,7,7) 4.8 & (14,4,4) 4.5 & (3,5,2) 3.6 & (1,2,5) 3.4 & (6,2,5) 2.9 \\
\bottomrule\end{tabular}
\caption{A3c: the five neurons per image as (map, row, column) and their activation.}\label{tab:a3c}
\end{table}

\begin{figure}[H]\centering
\includegraphics[width=0.80\linewidth]{A/A3c_patches_grid.png}
\infer{The patches sit on strokes, never on empty background, and the five patches of an image
cover different parts of the digit: the loop and the stem of the 6, the bar and the diagonal of the
7.  A $14\times14$ patch is half the image side, so each neuron sees a stroke segment or a junction,
not the whole digit.}
\caption{A3c: each image (rows) with the patch of each of its five most activated conv2 neurons
(columns).  An overlay with all five boxes per image is in \code{A3c\_patches\_overlay.png}.}\label{fig:a3c}
\end{figure}

% =====================================================================================
\section{Part B1: CNNs from Scratch on Caltech-101}\label{sec:b1}

\subsection{B1a -- Architectures and accuracy}
Table~\ref{tab:b1sizes} gives every layer's output size from Eq.~\eqref{eq:size} (each asserted
against the network).  All three share the first two convolution layers, and every convolution is
followed by ReLU.

\begin{table}[H]\centering\small\begin{tabular}{@{}lccc@{}}\toprule stage & Arch-1 & Arch-2 & Arch-3 \\ \midrule
input & $3\times224\times224$ & $3\times224\times224$ & $3\times224\times224$ \\
conv1 & $8\times54\times54$ & $8\times54\times54$ & $8\times54\times54$ \\
pool1 & $8\times26\times26$ & $8\times26\times26$ & $8\times26\times26$ \\
conv2 & $16\times22\times22$ & $16\times22\times22$ & $16\times22\times22$ \\
pool2 & $16\times10\times10$ & $16\times10\times10$ & $16\times10\times10$ \\
conv3 & -- & $32\times8\times8$ & $32\times8\times8$ \\
pool3 & -- & $32\times3\times3$ & -- \\
conv4 & -- & -- & $64\times6\times6$ \\
pool4 & -- & -- & $64\times2\times2$ \\
flatten & 1600 & 288 & 256 \\
fc1 & 128 & 128 & 128 \\
fc2 & 5 & 5 & 5 \\
\midrule parameters & 211,701 & 48,405 & 62,805 \\
last-conv receptive field & 51 px (stride 8) & 99 px (stride 16) & 131 px (stride 16) \\
\bottomrule\end{tabular}
\caption{B1 architectures: output size after each stage (channels $\times$ height $\times$ width),
parameter count and the receptive field of a last-convolution-layer neuron.
conv1 is 8 filters $11\times11$ stride 4 (as in AlexNet's first layer~\cite{alexnet}); conv2 16 filters $5\times5$; conv3 32 and conv4 64 filters
$3\times3$, all stride 1 padding 0; every pool is $3\times3$ stride 2; fc1 has 128 ReLU units.  Arch-3
has no pooling after conv3.  The flatten width is what drives the parameter count: Arch-1's
$1600\times128$ first dense layer holds most of its parameters.}\label{tab:b1sizes}
\end{table}

\begin{table}[H]\centering\small
\begin{subtable}{\linewidth}\centering\begin{tabular}{@{}lrrrrr@{}}\toprule
lr & best ep. & stop ep. & best val loss & val acc (\%) & train acc (\%) \\ \midrule
\rowcolor{selrow}0.0003 (edge) & 20 & 35 & 0.6853 & 74.0 (37/50) & 98.4 \\
0.001 & 9 & 24 & 0.8032 & 76.0 (38/50) & 94.8 \\
0.003 & 3 & 18 & 1.2013 & 56.0 (28/50) & 75.6 \\
0.01 & 3 & 18 & 1.4164 & 42.0 (21/50) & 42.8 \\
0.03 & 2 & 17 & 1.5797 & 22.0 (11/50) & 21.6 \\
0.1 & 17 & 32 & 1.6102 & 20.0 (10/50) & 20.0 \\
\bottomrule\end{tabular}\caption{Arch-1}\end{subtable}\\[4pt]
\begin{subtable}{\linewidth}\centering\begin{tabular}{@{}lrrrrr@{}}\toprule
lr & best ep. & stop ep. & best val loss & val acc (\%) & train acc (\%) \\ \midrule
\rowcolor{selrow}0.0003 (edge) & 31 & 46 & 0.5879 & 78.0 (39/50) & 97.2 \\
0.001 & 15 & 30 & 0.8429 & 74.0 (37/50) & 94.0 \\
0.003 & 4 & 19 & 1.3200 & 48.0 (24/50) & 54.8 \\
0.01 & 19 & 34 & 1.0891 & 46.0 (23/50) & 60.0 \\
0.03 & 14 & 29 & 1.6095 & 20.0 (10/50) & 20.0 \\
0.1 & 0 & 1 & 5.4059 & 20.0 (10/50) & 20.0 \\
\bottomrule\end{tabular}\caption{Arch-2}\end{subtable}\\[4pt]
\begin{subtable}{\linewidth}\centering\begin{tabular}{@{}lrrrrr@{}}\toprule
lr & best ep. & stop ep. & best val loss & val acc (\%) & train acc (\%) \\ \midrule
\rowcolor{selrow}0.0003 (edge) & 39 & 54 & 0.7256 & 74.0 (37/50) & 99.6 \\
0.001 & 13 & 28 & 0.7997 & 68.0 (34/50) & 94.8 \\
0.003 & 15 & 30 & 0.8885 & 70.0 (35/50) & 91.2 \\
0.01 & 10 & 25 & 1.2073 & 56.0 (28/50) & 63.2 \\
0.03 & 14 & 29 & 1.6095 & 20.0 (10/50) & 20.0 \\
0.1 & 0 & 1 & 3.3827 & 18.0 (9/50) & 18.4 \\
\bottomrule\end{tabular}\caption{Arch-3}\end{subtable}
\caption{B1 learning-rate grids.  Shaded rows are selected; ``edge'' marks the extra step tried
because the best grid rate was its lowest.  The rows at about 20\,\% (lr 0.03 and 0.1) collapsed:
the network ends predicting one class (loss $\approx \ln 5$), consistent with a step large enough to
switch the ReLUs off.  For Arch-2 and Arch-3 at 0.1 the loss diverged in the first epoch, so the
untrained weights (best epoch 0) are kept.}\label{tab:b1grid}
\end{table}

All three architectures selected the extra rate $0.0003$, so their best rate lies at or below the
lowest value tried.  The protocol allows one extension and was not changed after these results;
Section~\ref{sec:limits} discusses what this means for the comparison.

\begin{table}[H]\centering\small\begin{tabular}{@{}lrrrrr@{}}\toprule model & params & lr & train acc & val acc & val loss \\ \midrule
Arch-1 & 211,701 & 0.0003 & 98.4 (246/250) & 74.0 (37/50) & 0.685 \\
Arch-2 & 48,405 & 0.0003 & 97.2 (243/250) & 78.0 (39/50) & 0.588 \\
Arch-3 & 62,805 & 0.0003 & 99.6 (249/250) & 74.0 (37/50) & 0.726 \\
\bottomrule\end{tabular}
\caption{B1a: training and validation accuracy at the selected learning rate.}\label{tab:b1acc}
\end{table}

\begin{figure}[H]\centering
\begin{subfigure}{0.32\linewidth}\includegraphics[width=\linewidth]{B1/B1a_arch1_cm_train.png}\caption{Arch-1 train}\end{subfigure}\hfill
\begin{subfigure}{0.32\linewidth}\includegraphics[width=\linewidth]{B1/B1a_arch2_cm_train.png}\caption{Arch-2 train}\end{subfigure}\hfill
\begin{subfigure}{0.32\linewidth}\includegraphics[width=\linewidth]{B1/B1a_arch3_cm_train.png}\caption{Arch-3 train}\end{subfigure}\\[4pt]
\begin{subfigure}{0.32\linewidth}\includegraphics[width=\linewidth]{B1/B1a_arch1_cm_val.png}\caption{Arch-1 validation}\end{subfigure}\hfill
\begin{subfigure}{0.32\linewidth}\includegraphics[width=\linewidth]{B1/B1a_arch2_cm_val.png}\caption{Arch-2 validation}\end{subfigure}\hfill
\begin{subfigure}{0.32\linewidth}\includegraphics[width=\linewidth]{B1/B1a_arch3_cm_val.png}\caption{Arch-3 validation}\end{subfigure}
\caption{B1a confusion matrices.}\label{fig:b1cm}
\end{figure}

\begin{figure}[H]\centering
\includegraphics[width=0.85\linewidth]{B1/B1a_arch2_curves.png}
\infer{The gap between training and validation opens within the first ten epochs and keeps
widening: training loss falls steadily while validation loss levels off.  With 250 training images
and no augmentation or regularisation in the specified networks, early stopping is the only thing
that limits overfitting, and the three networks reach 97.2--99.6\,\% on
the training set (Table~\ref{tab:b1acc}).}
\caption{Training curves of the selected Arch-2 run (the other two architectures, in the code
output, show the same pattern).}\label{fig:b1curves}
\end{figure}

\subsection{B1b -- Best architecture on test}
By validation accuracy (ties broken by validation loss, then fewer parameters), the best
architecture is Arch-2 (5-layer): 39/50 correct against
37 for Arch-1 and 37 for Arch-3.  On the test set it
scores 68.0\,\% (68/100).  Per-class test recall is
bonsai 80\,\%, buddha 55\,\%, car\_side
85\,\%, scorpion 50\,\% and starfish
70\,\%.  Arch-2 is the one the protocol selects, but its margin is 2 validation images from a
single run of each architecture; the seed study below shows that the three perform similarly.

\begin{figure}[H]\centering
\includegraphics[width=0.45\linewidth]{B1/B1b_best_cm_test.png}
\infer{The test accuracy is below the validation accuracy that selected the model; that drop is
expected, because the validation set was used to choose the epoch, the learning rate and the
architecture (Section~\ref{sec:limits}).  car\_side, the one entirely greyscale class, has the
highest recall, which may partly reflect that the absence of colour alone identifies it.  The most confused pair is scorpion and starfish
(8 of the 32 errors, both directions together), and the largest single
cell is scorpion predicted as bonsai (5).}
\caption{B1b: test confusion matrix of Arch-2 (5-layer).}\label{fig:b1test}
\end{figure}

\paragraph{Additional experiment: stability over seeds.}\label{sec:seeds}
Because the selection rests on 50 validation images and one training run, the whole B1 protocol
(learning-rate grid, early stopping, selection rule) was rerun unchanged for five seeds
(\code{seeds\_b1.py}); a seed changes only the initial weights and the order of the training
images, and seed 11 is the run reported above.  Within each seed the architecture is chosen on
validation as before; the test accuracy of all three is recorded afterwards for this analysis only.

\begin{table}[H]\centering\small\begin{tabular}{@{}lrrrrrrl@{}}\toprule
 & \multicolumn{2}{c}{Arch-1} & \multicolumn{2}{c}{Arch-2} & \multicolumn{2}{c}{Arch-3} & \\
\cmidrule(lr){2-3}\cmidrule(lr){4-5}\cmidrule(lr){6-7}
seed & val & test & val & test & val & test & selected \\ \midrule
11 & 74 & 68 & 78 & 68 & 74 & 69 & Arch-2 \\
12 & 80 & 72 & 72 & 77 & 80 & 70 & Arch-1 \\
13 & 72 & 71 & 84 & 73 & 86 & 73 & Arch-3 \\
14 & 74 & 70 & 74 & 68 & 66 & 69 & Arch-2 \\
15 & 76 & 69 & 80 & 72 & 74 & 73 & Arch-2 \\ \midrule
mean $\pm$ sd & 75.2\,$\pm$\,3.0 & 70.0\,$\pm$\,1.6 & 77.6\,$\pm$\,4.8 & 71.6\,$\pm$\,3.8 & 76.0\,$\pm$\,7.5 & 70.8\,$\pm$\,2.0 & \\
\bottomrule\end{tabular}
\caption{B1 over five seeds: validation and test accuracy (\%) of each architecture at its selected
learning rate, and the architecture the validation rule selects.  sd is the sample standard
deviation across seeds.  Ties in validation accuracy (seeds 12 and 14) are broken by validation
loss.}\label{tab:seeds}
\end{table}

The validation rule picks Arch-2 three times and Arch-1 and Arch-3 once each, and one
architecture's validation accuracy alone moves by up to 20 points between seeds (Arch-3, 66--86\,\%).
The mean test accuracies, 70.0, 71.6 and 70.8\,\%, differ by at most 1.6 points, no more than
the seed-to-seed standard deviations (1.6--3.8 points), so the data give no evidence that any architecture is better (paired by seed, Arch-2 minus
Arch-1 is $0$, $+5$, $+2$, $-2$ and $+3$ points).  The network chosen on validation averages
70.6\,\% on test (68--73\,\%); the 68.0\,\% of the reported run is at the low end of Arch-2's range.

\subsection{B1c -- Feature maps of one training image}
For one bonsai training image, all 8 conv1 maps are shown, and for each later convolution layer
the 8 maps with the highest mean activation: conv2 maps 5, 3, 8, 9, 14, 7, 12, 10 and conv3 maps
9, 13, 29, 11, 24, 12, 25, 22 (in that order, also in the panel titles).  Maps are after ReLU and before
pooling, on one colour scale per layer.

\begin{figure}[H]\centering
\includegraphics[width=0.78\linewidth]{B1/B1c_conv1_maps.png}\\[2pt]
\includegraphics[width=0.78\linewidth]{B1/B1c_conv2_maps.png}\\[2pt]
\includegraphics[width=0.78\linewidth]{B1/B1c_conv3_maps.png}
\infer{conv1 maps are sharp and image-like: they separate the plant from the plain background, and
individual maps pick out the pot rim, the foliage or the trunk texture.  Moving up, the maps get
coarser ($54\times54$, $22\times22$, $8\times8$) and denser: over all maps of each layer, 66\,\% of
the conv1 values are exactly zero, 60\,\% of conv2 and 39\,\% of conv3.  (The maps shown for conv2
and conv3 are the most active ones, so they are denser than their layer: 44\,\% and 5\,\% zeros.)  The deeper maps no longer outline the object; they respond over
large regions, and the strongest conv3 map is active almost everywhere.}
\caption{B1c: conv1 (top, $8\times54\times54$), conv2 (middle, $16\times22\times22$) and conv3 (bottom,
$32\times8\times8$) feature maps of Arch-2 for the bonsai image.}\label{fig:b1c}
\end{figure}

\subsection{B1d -- Patches that maximally activate the last convolution layer}
For one training image per class, the single most activated neuron of conv3 (Arch-2's last
convolution layer, $32\times8\times8$) is traced back with Eq.~\eqref{eq:trace}.  Its receptive
field is 99$\times$99 pixels, and neighbouring neurons are 16 pixels
apart.  As in A3c, the neuron is chosen per image, so each patch shows what fired most on that
image, not what one fixed neuron prefers; map 9 is the winner for four of the five images.

\begin{figure}[H]\centering
\includegraphics[width=0.38\linewidth]{B1/B1d_patches.png}
\infer{Every patch lies on the object: the trunk and lower foliage of the bonsai, the shoulder of
the buddha, the body and wheels of the car, the claws of the scorpion and a starfish.  The same
feature map (map 9) is the most active unit for four of the five classes, so it is not specific
to one class.  The starfish neuron fires far more weakly (3.44, against 7.7--9.2 for the others)
and from a different map, possibly because the image is low in contrast (small sand-coloured
starfish on sand).}
\caption{B1d: each image with the patch of its most activated conv3 neuron (left) and the patch
itself (right).}\label{fig:b1d}
\end{figure}

\begin{table}[H]\centering\small\begin{tabular}{@{}lrrrl@{}}\toprule image & map & (row, col) & activation & input patch \\ \midrule
bonsai & 9 & (4, 5) & 7.71 & rows 49--147, cols 65--163 \\
buddha & 9 & (5, 7) & 8.80 & rows 65--163, cols 97--195 \\
car\_side & 9 & (5, 4) & 9.05 & rows 65--163, cols 49--147 \\
scorpion & 9 & (8, 6) & 9.19 & rows 113--211, cols 81--179 \\
starfish & 25 & (6, 2) & 3.44 & rows 81--179, cols 17--115 \\
\bottomrule\end{tabular}
\caption{B1d: the most activated conv3 neuron per image and its input patch (1-based pixel
rows and columns).}\label{tab:b1d}
\end{table}

% =====================================================================================
\section{Part B2: Transfer Learning with VGG19}\label{sec:b2}

\subsection{Model}
\code{keras.applications.VGG19}~\cite{vgg} with ImageNet weights is used with its full classifier.  Only the
final 1000-way layer is replaced, by a 5-way softmax layer, and only that layer is trained
(20,485 of 139,590,725 parameters); this follows the lecture notes' advice to retrain
only the softmax layer when the dataset is small.  Since everything below the new layer is frozen,
deterministic and free of dropout, its 4096-dimensional input (the \code{fc2} output) is computed
once per image and the layer is trained on those values.  This is the same training as running
the frozen network at every step, and the code checks that the assembled model gives identical
predictions.

\begin{table}[H]\centering\small\begin{tabular}{@{}lrrrrr@{}}\toprule
lr & best ep. & stop ep. & best val loss & val acc (\%) & train acc (\%) \\ \midrule
0.001 & 1 & 16 & 0.0296 & 98.0 (49/50) & 98.0 \\
0.003 & 500 & 500 & 0.0081 & 100.0 (50/50) & 100.0 \\
0.01 & 1 & 16 & 0.0163 & 98.0 (49/50) & 100.0 \\
\rowcolor{selrow}0.03 & 4 & 19 & $2.4\times10^{-8}$ & 100.0 (50/50) & 100.0 \\
0.1 & 4 & 19 & 1.0027 & 96.0 (48/50) & 100.0 \\
\bottomrule\end{tabular}
\caption{B2 learning-rate grid for the retrained layer.}\label{tab:b2grid}
\end{table}

The selected rate 0.03 reaches a validation loss of $2.4\times10^{-8}$ at epoch
4: every validation image is classified with probability close to 1.  At this
point the grid can no longer tell the rates apart: every rate from 0.001 to 0.03 reaches at least
98\,\% validation accuracy, and the run at 0.003 stopped only at the 500-epoch limit
because its validation loss kept falling by tiny amounts.

\subsection{B2a -- Accuracy and comparison with B1b}
The VGG19 model classifies 250/250 training,
50/50 validation and 100/100 test
images correctly.

\begin{figure}[H]\centering
\begin{subfigure}{0.32\linewidth}\includegraphics[width=\linewidth]{B2/B2a_cm_train.png}\caption{train}\end{subfigure}\hfill
\begin{subfigure}{0.32\linewidth}\includegraphics[width=\linewidth]{B2/B2a_cm_val.png}\caption{validation}\end{subfigure}\hfill
\begin{subfigure}{0.32\linewidth}\includegraphics[width=\linewidth]{B2/B2a_cm_test.png}\caption{test}\end{subfigure}
\caption{B2a confusion matrices of VGG19 with the retrained 5-class layer.}\label{fig:b2cm}
\end{figure}

\begin{table}[H]\centering\small\begin{tabular}{@{}lrrrr@{}}\toprule model & trainable params & train & val & test \\ \midrule
Arch-2 (5-layer) (from scratch) & 48,405 & 97.2 & 78.0 & 68.0 \\
VGG19 ImageNet, head retrained & 20,485 & 100.0 & 100.0 & 100.0 \\
\bottomrule\end{tabular}
\caption{B2a: VGG19 against the best network trained from scratch, on the same images.}\label{tab:b2cmp}
\end{table}

\begin{figure}[H]\centering
\includegraphics[width=0.62\linewidth]{B2/B2a_comparison.png}
\infer{The gap is 100.0 against 68.0\,\% on the same 100 test images,
far beyond what the sampling noise of 100 images can explain.  It is a comparison of two ways of
obtaining features, not of two architectures: VGG19's convolutional features were learned from
over a million ImageNet images and are fixed, while the B1 networks learn theirs from 250 images;
Arch-2 trains 48,405 weights against VGG19's 20,485.  The task is also close to ImageNet's own:
starfish and scorpion are ImageNet classes, and it has car and pot classes, so VGG19's features
already separate these five classes.}
\caption{B2a: train, validation and test accuracy of the best B1 network and of VGG19.}\label{fig:b2cmp}
\end{figure}

\subsection{B2b -- Patches that maximally activate the last convolution layer}
The same five images as B1d are passed through VGG19 and the most activated neuron of
\code{block5\_conv4} ($14\times14\times512$) is traced back.  Its receptive field is
252~pixels, larger than the $224$-pixel image, so every traced box is clipped by the image
border and covers most of the image.  To show where the neuron actually looks, the figure adds the
magnitude of the gradient of the neuron with respect to each input pixel; pixels above
$10^{-3}$ of the maximum gradient cover 43--57\,\% of the
image.

\begin{figure}[H]\centering
\includegraphics[width=0.55\linewidth]{B2/B2b_patches.png}
\infer{For a neuron this deep the receptive-field box carries almost no information, since it is
close to the whole image.  The gradient map is what localises the neuron: its strongest values
sit on the object (the trunk and pot of the bonsai, the left side of the buddha image with its arm
and the relief beside it, the car body, the scorpion's body, the nearest starfish), even though the
non-zero gradient still spreads over about half of the image.}
\caption{B2b: image with the clipped receptive-field box (left), the box contents (middle) and the
neuron's input-gradient magnitude (right).}\label{fig:b2b}
\end{figure}

\begin{table}[H]\centering\small\begin{tabular}{@{}lrrrl@{}}\toprule image & map & (row, col) & activation & input box (clipped) \\ \midrule
bonsai & 450 & (11, 8) & 110 & rows 43--224, cols 1--224 \\
buddha & 237 & (10, 3) & 131 & rows 27--224, cols 1--166 \\
car\_side & 451 & (8, 4) & 137 & rows 1--224, cols 1--182 \\
scorpion & 105 & (6, 10) & 133 & rows 1--214, cols 27--224 \\
starfish & 499 & (10, 10) & 58 & rows 27--224, cols 27--224 \\
\bottomrule\end{tabular}
\caption{B2b: the most activated \code{block5\_conv4} neuron per image and its (clipped) input box.}\label{tab:b2b}
\end{table}

\subsection{B2c -- Guided backpropagation}
For the B1c bonsai image, five \code{block5\_conv4} neurons with non-zero activation are chosen:
feature maps in order of peak activation, each at its peak, skipping any within 3 positions of a
neuron already chosen (the strongest maps all peak around the same object part).  Keeping only the
chosen neuron and backpropagating to the input, guided backpropagation~\cite{springenberg} passes a gradient through a
ReLU only where both the forward input and the incoming gradient are positive,
\begin{equation}\label{eq:guided}
R^{(l)} = \mathbb 1\!\left[x^{(l)} > 0\right]\,\mathbb 1\!\left[R^{(l+1)} > 0\right]\,R^{(l+1)} ,
\end{equation}
which removes the negative influences that make the plain gradient noisy.  The convolution weights
are ImageNet's (they are unchanged by B2a, which trains only the last layer).

\begin{table}[H]\centering\small\begin{tabular}{@{}rrrr@{}}\toprule neuron & (row, col) in $14\times14$ & feature map & activation \\ \midrule
1 & (11, 8) & 450 & 109.7 \\
2 & (1, 1) & 120 & 61.1 \\
3 & (12, 4) & 16 & 59.0 \\
4 & (8, 3) & 149 & 53.4 \\
5 & (14, 1) & 230 & 49.7 \\
\bottomrule\end{tabular}
\caption{B2c: the five neurons, as (row, column) in the $14\times14$ grid of a feature map.}\label{tab:b2c}
\end{table}

\begin{figure}[H]\centering
\includegraphics[width=\linewidth]{B2/B2c_guided_backprop.png}
\infer{Each neuron is driven by a small, specific part of the image: the pot and the base of the
trunk, one of the light-green cones, the left rim of the pot, a cone on the left branch, the pot's lower
left corner.  The neuron at grid position (1,\,1) responds to the top-left corner of the frame and
the cone nearest to it.  The guided images show those parts as sharp edges and colour, whereas the
plain gradients (bottom row) cover the same regions but as noise.  Even at the last convolution
layer, a single neuron encodes a part, not the object.  Guided backpropagation is used here as a
visualisation: its sharp, image-like detail partly comes from the ReLU masking itself, so it shows
where a neuron looks, not a faithful measure of what it uses.}
\caption{B2c: guided backpropagation (top) and the plain gradient (bottom) for the five neurons.}\label{fig:b2c}
\end{figure}

\subsection{B2d -- Grad-CAM}
Grad-CAM~\cite{gradcam}, following the Keras example~\cite{keras}, weights each \code{block5\_conv4} feature map $A^k$ by the
average gradient of the class score $y^c$ (before softmax) over its positions and keeps the positive
part of the weighted sum:
\begin{equation}\label{eq:gradcam}
\alpha^c_k = \frac{1}{14^2}\sum_{i,j}\frac{\partial y^c}{\partial A^k_{ij}},\qquad
L^c = \mathrm{ReLU}\Big(\sum_k \alpha^c_k A^k\Big).
\end{equation}
For each of the five B1d images a map is computed for each of the five classes (25 maps), using
the VGG19 model with its retrained 5-class layer.  Each map is scaled to its own maximum, with the
raw maximum printed above it.

\begin{figure}[H]\centering
\includegraphics[width=0.9\linewidth]{B2/B2d_gradcam.png}
\infer{For every image the true class has the largest raw Grad-CAM maximum, and its map lies on
the class's defining part: the wheels of the car, the head of the buddha, the body of the scorpion,
the largest starfish.  For bonsai, the evidence sits on the pot and soil more than on the tree,
which suggests the model recognises a bonsai partly by its pot.  Maps for the wrong classes are
weaker and sometimes fall on parts that resemble those classes (the starfish map of the buddha falls on the
torso, the scorpion map of the starfish image on an arm of the large starfish), while others land on
the object itself (the bonsai map of the starfish image, the scorpion map of the car).  Because every map is
scaled to its own maximum, only the printed raw maxima can be compared across panels.}
\caption{B2d Grad-CAM: each image (rows) against each target class (columns); the true class is
framed.}\label{fig:b2d}
\end{figure}

\begin{table}[H]\centering\small\begin{tabular}{@{}llr@{}}\toprule image (true) & predicted & $p(\text{true})$ \\ \midrule
bonsai & bonsai & 1.000 \\
buddha & buddha & 1.000 \\
car\_side & car\_side & 1.000 \\
scorpion & scorpion & 1.000 \\
starfish & starfish & 1.000 \\
\bottomrule\end{tabular}
\caption{B2d: prediction for each Grad-CAM image.}\label{tab:b2d}
\end{table}

% =====================================================================================
\section{Limitations}\label{sec:limits}
\begin{itemize}[leftmargin=*,itemsep=2pt]
\item \textbf{Validation numbers are optimistic; test numbers are not.}  The 50 validation images
  chose each model's stopping epoch and learning rate and, in B1b, the architecture.  Reported
  validation accuracies are therefore biased upwards, which is visible in B1b
  (78.0\,\% validation, 68.0\,\% test).  The test set was used only
  after every choice was made, so the test accuracies are unbiased estimates.
\item \textbf{The B1 architecture ranking is not supported.}  In the reported run Arch-2 was
  chosen by 2 validation images (39 vs 37/50), and at learning rate 0.001 Arch-1 would have won
  (Table~\ref{tab:b1grid}).  Over five seeds (Table~\ref{tab:seeds}) the selection changes and the
  mean test accuracies differ by at most 1.6 points, no more than the seed-to-seed spread, so the three architectures
  perform similarly on this data.  In the reported run all three chose the extra rate 0.0003, so
  their best rate may be lower still; over the five seeds 0.0003 was chosen in 8 of 15 runs.
\item \textbf{car\_side is entirely greyscale.}  Its recognition can rest partly on the absence of
  colour, which may contribute to its high recall in B1b.
\item \textbf{The VGG19 learning rate was chosen at a ceiling.}  With all rates near 100\,\%
  validation accuracy, the selected rate won on a validation loss of order $10^{-8}$ that lasted one
  epoch (epoch 4; it was back above $10^{-3}$ at epoch 5); any rate from 0.001 to 0.03 would give the same conclusion in B2a.
\item \textbf{Receptive fields beyond the image.}  For \code{block5\_conv4} the traced box is
  clipped to almost the whole image, so B2b's localisation comes from the gradient map, not from
  the box.
\end{itemize}

% =====================================================================================
\section{Conclusion}\label{sec:conclusion}
Building the network from its parts made each part's role visible: a random filter already
responds to edges, ReLU and pooling turn those responses into sparse, lower-resolution maps, and
training sharpens the filters without discarding their initial structure, reaching
99.42\,\% on the MNIST subset.  On Caltech-101 with 250 training images, three CNNs
trained from scratch all overfit heavily, and the one selected on validation reaches 68.0\,\% on test
(70--72\,\% on average for every architecture over five seeds), while
ImageNet features with a retrained final layer reach 100.0\,\%.  On this data, the
features a network starts with matter far more than which of the three specified architectures it
uses.  The
visualisations agree with this from the inside: VGG19's deepest neurons respond to object parts,
and its class evidence sits on the parts that define each class.

% =====================================================================================
\begin{thebibliography}{9}
\bibitem{lecun} Y.~LeCun, L.~Bottou, Y.~Bengio and P.~Haffner, ``Gradient-based learning applied
  to document recognition,'' \emph{Proceedings of the IEEE}, 86(11):2278--2324, 1998.
\bibitem{alexnet} A.~Krizhevsky, I.~Sutskever and G.~E.~Hinton, ``ImageNet classification with
  deep convolutional neural networks,'' in \emph{NeurIPS}, 2012.
\bibitem{vgg} K.~Simonyan and A.~Zisserman, ``Very deep convolutional networks for large-scale
  image recognition,'' arXiv:1409.1556, 2014.
\bibitem{he} K.~He, X.~Zhang, S.~Ren and J.~Sun, ``Delving deep into rectifiers: surpassing
  human-level performance on ImageNet classification,'' in \emph{ICCV}, 2015.
\bibitem{glorot} X.~Glorot and Y.~Bengio, ``Understanding the difficulty of training deep
  feedforward neural networks,'' in \emph{AISTATS}, 2010.
\bibitem{zeiler} M.~D.~Zeiler and R.~Fergus, ``Visualizing and understanding convolutional
  networks,'' in \emph{ECCV}, 2014.
\bibitem{springenberg} J.~T.~Springenberg, A.~Dosovitskiy, T.~Brox and M.~Riedmiller, ``Striving
  for simplicity: the all convolutional net,'' arXiv:1412.6806, 2014.
\bibitem{gradcam} R.~R.~Selvaraju, M.~Cogswell, A.~Das, R.~Vedantam, D.~Parikh and D.~Batra,
  ``Grad-CAM: visual explanations from deep networks via gradient-based localization,'' in
  \emph{ICCV}, 2017.
\bibitem{keras} F.~Chollet, ``Grad-CAM class activation visualization,''
  \url{https://keras.io/examples/vision/grad_cam/}.
\end{thebibliography}

\end{document}
